% problem_id: S001-spaces-morphisms-lemma-refined-valuative-criterion-universally-closed
% source_id: S001 (Stacks Project)
% source_file: spaces-morphisms.tex
% statement_locator: spaces-morphisms.tex lines 8778--8812
% proof_locator: spaces-morphisms.tex lines 8814--8943
% env: lemma
% id_note: label 跨文件重名 ⇒ id 用文件限定形式
% pairing: tight (verified)
% status: complete (statement + author proof verbatim + reason + locators)
%
\begin{lemma}
\label{lemma-refined-valuative-criterion-universally-closed}
Let $S$ be a scheme. Let $f : X \to Y$ and $h : U \to X$ be morphisms of
algebraic spaces over $S$. If
\begin{enumerate}
\item $f$ and $h$ are quasi-compact,
\item $|h|(|U|)$ is dense in $|X|$, and
\end{enumerate}
given any commutative solid diagram
$$
\xymatrix{
\Spec(K) \ar[r] \ar[d] & U \ar[r] & X \ar[d] \\
\Spec(A) \ar[rr] \ar@{-->}[rru] & & Y
}
$$
where $A$ is a valuation ring with field of fractions $K$
\begin{enumerate}
\item[(3)] there exists at most one dotted arrow making the diagram
commute, and
\item[(4)] there exists an extension $K'/K$ of fields, a
valuation ring $A' \subset K'$ dominating $A$ and a morphism
$\Spec(A') \to X$ such that the following diagram commutes
$$
\xymatrix{
\Spec(K') \ar[r] \ar[d] & \Spec(K) \ar[r] & U \ar[r] & X \ar[d] \\
\Spec(A') \ar[r] \ar[rrru] & \Spec(A) \ar[rr] & & Y
}
$$
\end{enumerate}
then $f$ is universally closed. If moreover
\begin{enumerate}
\item[(5)] $f$ is quasi-separated
\end{enumerate}
then $f$ is separated and universally closed.
\end{lemma}

\begin{proof}
Assume (1), (2), (3), and (4).
We will verify the existence part of the valuative criterion for $f$
which will imply $f$ is universally closed by
Lemma \ref{lemma-quasi-compact-existence-universally-closed}.
To do this, consider a commutative diagram
\begin{equation}
\label{equation-start-with}
\vcenter{
\xymatrix{
\Spec(K) \ar[r] \ar[d] & X \ar[d] \\
\Spec(A) \ar[r] & Y
}
}
\end{equation}
where $A$ is a valuation ring and $K$ is the fraction field of $A$.
Note that since valuation rings and fields are reduced, we may
replace $U$, $X$, and $S$ by their respective reductions by
Properties of Spaces, Lemma \ref{spaces-properties-lemma-map-into-reduction}.
In this case the assumption that $h(U)$ is dense means that
the scheme theoretic image of $h : U \to X$ is $X$, see
Lemma \ref{lemma-scheme-theoretic-image-reduced}.

\medskip\noindent
Reduction to the case $Y$ affine. Choose an \'etale morphism
$\Spec(R) \to Y$ such that the closed point of $\Spec(A)$ maps
to an element of $\Im(|\Spec(R)| \to |Y|)$. By
Lemma \ref{lemma-lift-valuation-ring-through-flat-morphism}
we can find a local ring map $A \to A'$ of valuation rings
and a morphism $\Spec(A') \to \Spec(R)$ fitting into a commutative
diagram
$$
\xymatrix{
\Spec(A') \ar[r] \ar[d] & \Spec(R) \ar[d] \\
\Spec(A) \ar[r] & Y
}
$$
Since in Definition \ref{definition-valuative-criterion}
we allow for extensions of valuation rings
it is clear that we may replace $A$ by $A'$, $Y$ by $\Spec(R)$,
$X$ by $X \times_Y \Spec(R)$ and $U$ by $U \times_Y \Spec(R)$.

\medskip\noindent
From now on we assume that $Y = \Spec(R)$ is an affine scheme.
Let $\Spec(B) \to X$ be an \'etale morphism from an affine scheme
such that the morphism $\Spec(K) \to X$ is in the image of
$|\Spec(B)| \to |X|$. Since we may replace $K$ by an extension
$K' \supset K$ and $A$ by a valuation ring $A' \subset K'$
dominating $A$ (which exists by
Algebra, Lemma \ref{algebra-lemma-dominate}),
we may assume the morphism $\Spec(K) \to X$ factors through $\Spec(B)$
(by definition of $|X|$). In other words, we may think of $K$ as a $B$-algebra.
Choose a polynomial algebra $P$ over $B$ and a $B$-algebra surjection
$P \to K$. Then $\Spec(P) \to X$ is flat as a composition
$\Spec(P) \to \Spec(B) \to X$. Hence the scheme theoretic image
of the morphism $U \times_X \Spec(P) \to \Spec(P)$ is $\Spec(P)$ by
Lemma \ref{lemma-flat-base-change-scheme-theoretic-image}.
By Lemma \ref{lemma-reach-points-scheme-theoretic-image}
we can find a commutative diagram
$$
\xymatrix{
\Spec(K') \ar[r] \ar[d] & U \times_X \Spec(P) \ar[d] \\
\Spec(A') \ar[r] & \Spec(P)
}
$$
where $A'$ is a valuation ring and $K'$ is the fraction field of $A'$
such that the closed point of $\Spec(A')$ maps to
$\Spec(K) \subset \Spec(P)$. In other words, there is a $B$-algebra map
$\varphi : K \to A'/\mathfrak m_{A'}$. Choose a valuation ring
$A'' \subset A'/\mathfrak m_{A'}$ dominating $\varphi(A)$ with
field of fractions $K'' = A'/\mathfrak m_{A'}$
(Algebra, Lemma \ref{algebra-lemma-dominate}). We set
$$
C = \{\lambda \in A' \mid \lambda \bmod \mathfrak m_{A'} \in A''\}.
$$
which is a valuation ring by
Algebra, Lemma \ref{algebra-lemma-stack-valuation-rings}.
As $C$ is an $R$-algebra with fraction field $K'$, we obtain a solid
commutative diagram
$$
\xymatrix{
\Spec(K'_1) \ar@{-->}[r] \ar@{-->}[d] &
\Spec(K') \ar[r] \ar[d] & U \ar[r] & X \ar[d] \\
\Spec(C_1) \ar@{-->}[r] \ar@{-->}[rrru] & \Spec(C) \ar[rr] & & Y
}
$$
as in the statement of the lemma. Thus assumption (4) produces
$C \to C_1$ and the dotted arrows making the diagram commute.
Let $A_1' = (C_1)_\mathfrak p$ be the localization of $C_1$
at a prime $\mathfrak p \subset C_1$ lying over
$\mathfrak m_{A'} \subset C$. Since $C \to C_1$ is flat by
More on Algebra, Lemma
\ref{more-algebra-lemma-valuation-ring-torsion-free-flat}
such a prime $\mathfrak p$ exists by
Algebra, Lemmas \ref{algebra-lemma-local-flat-ff} and
\ref{algebra-lemma-ff-rings}.
Note that $A'$ is the localization of $C$ at $\mathfrak m_{A'}$
and that $A'_1$ is a valuation ring
(Algebra, Lemma \ref{algebra-lemma-make-valuation-rings}).
In other words, $A' \to A'_1$ is a local ring map of valuation
rings. Assumption (3) implies
$$
\xymatrix{
\Spec(A'_1) \ar[r] \ar[d] & \Spec(C_1) \ar[r] & X \\
\Spec(A') \ar[r] & \Spec(P) \ar[r] & \Spec(B) \ar[u]
}
$$
commutes. Hence the restriction of the morphism $\Spec(C_1) \to X$
to $\Spec(C_1/\mathfrak p)$ restricts to the composition
$$
\Spec(\kappa(\mathfrak p)) \to
\Spec(A'/\mathfrak m_{A'}) = \Spec(K'') \to
\Spec(K) \to X
$$
on the generic point of $\Spec(C_1/\mathfrak p)$. Moreover,
$C_1/\mathfrak p$ is a valuation ring
(Algebra, Lemma \ref{algebra-lemma-make-valuation-rings})
dominating $A''$ which dominates $A$.
Thus the morphism $\Spec(C_1/\mathfrak p) \to X$ witnesses the
existence part of the valuative criterion for the diagram
(\ref{equation-start-with}) as desired.

\medskip\noindent
Next, suppose that (5) is satisfied as well, i.e., the morphism
$\Delta : X \to X \times_S X$ is quasi-compact. In this case
assumptions (1) -- (4) hold for $h$ and $\Delta$. Hence the first
part of the proof shows that $\Delta$ is universally closed.
By Lemma \ref{lemma-separated-diagonal-proper} we conclude that
$f$ is separated.
\end{proof}
